\chapter*{Abstract}
\addcontentsline{toc}{chapter}{Abstract}

In Bulgaria, over 6,700 children have speech impairment \cite{unicef_bulgaria_aac}. They are mainly helped by speech therapists, but weekly or biweekly meetings are often insufficient. Practice at home can be scarce and difficult, and without parental supervision, which often does not leave enough time, children are left without the necessary help \cite{home_practice_pubmed}. Lack of attention at home, or more in-depth analysis or supervision by a speech therapist, can further deepen speech impairment in children, which could result in problems speaking at a later age. To address this limitation, in our Bachelor Project, we developed and tested a software application based on self-study and gamification concepts to empower children with such an impairment to complete their exercises independently, while still under the oversight of their therapist and parents. Our solution showed promising results, and our evaluation includes expert feedback from a speech therapist, indicating the potential to improve motivation among kids to undertake exercises independently and to optimize the therapist's efficiency by delegating exercises and tracking children's results. This application showed promising results contingent on future improvements in the implementation of gamification and the extent to which self-study can be sustained without direct observation.
